\section{Citations and References}
\label{sec:citations}
Citations make clear which ideas, methods, results, and resources come
from existing work and which are the author's own. They allow readers to
verify claims, locate the work on which the report builds, and distinguish
established knowledge from new contributions. This section explains when
citations are required, how to manage references with \LaTeX{}, and how to
cite software, datasets, and AI tools appropriately.

\subsection{Why and when to cite}
Cite a source whenever you rely on a fact, method, figure, or idea that
is not your own or not common knowledge in the field. This includes:
background facts (e.g.\ the original definition of an algorithm
\cite{dijkstra1959}), tools and datasets you used, and prior work you
compare against. A rule of thumb: if you could not have written the
sentence without having read that source, cite it. Foundational results
you build directly on should also be cited, even if they are famous ---
for example, referencing Shannon's original formulation when discussing
channel capacity \cite{shannon1948}.

\begin{pitfallbox}
  Citing a source is not the same as quoting it. Paraphrase in your own
  words and cite; use direct quotation marks only for the rare case where
  the exact wording matters (e.g.\ a definition from a standard).
  Copy-pasting text and adding a citation at the end of the paragraph is
  still plagiarism if the wording is not your own.
\end{pitfallbox}

\subsection{Using \texttt{natbib} and BibTeX}
This template uses \texttt{natbib} with the \texttt{ieeetr} bibliography
style, a numeric style common in CS/engineering venues. Entries live in
\texttt{references.bib}; a minimal entry looks like:

\begin{lstlisting}[language=TeX, caption={A BibTeX entry.}]
@inproceedings{smith2023scheduling,
  author    = {Smith, Jane and Doe, John},
  title     = {Load-Aware Scheduling for Microservices},
  booktitle = {Proc. of ExampleConf},
  year      = {2023},
  pages     = {1--12},
}
\end{lstlisting}

and is cited in the text with \verb|\cite{smith2023scheduling}|, which
renders as a bracketed number with the \texttt{ieeetr} style used here,
e.g.\ the classic paper on machine intelligence \cite{turing1950}. (If
you use an author--year style such as \texttt{plainnat} or
\texttt{biblatex}'s \texttt{authoryear} instead, \texttt{natbib}'s
\verb|\citet{...}| lets you use the author name as a noun in the
sentence, e.g.\ ``Turing \verb|\citet{turing1950}| asked whether
machines can think.'') Every entry you cite must appear in
\texttt{references.bib},
and (with numeric styles) every entry in the bibliography should be
cited somewhere --- do not pad the reference list.

\begin{tipbox}
  If your own \LaTeX{} install has \texttt{biblatex}/\texttt{biber}
  available (true on Overleaf and most modern TeX Live installs), you can
  swap this template's \texttt{natbib} block in \texttt{main.tex} for
  \verb|\usepackage[backend=biber,style=ieee]{biblatex}| plus
  \verb|\addbibresource{references.bib}| and
  \verb|\printbibliography|; \texttt{biblatex} offers more citation
  styles and better internationalisation. Both approaches read the same
  \texttt{references.bib} file.
\end{tipbox}

\subsection{Citing software, datasets, and AI tools}
Cite the tools and datasets that materially shaped your results, e.g.\ a
specific library version, dataset, or pretrained model, the same way you
would cite a paper --- most have a recommended citation in their
documentation or a Zenodo/DOI record. If you used a generative AI tool
(e.g.\ for code suggestions or editing), follow your course's or
institution's disclosure policy; when in doubt, state in a short
paragraph (e.g.\ at the end of the introduction or in an
acknowledgements section) what the tool was used for and confirm you
verified the result yourself.

\begin{tipbox}
  Keep an eye on citation style consistency: IEEE numeric
  (\texttt{[1]}), ACM, or author--year (APA/Chicago) --- pick the one your
  course requires and use it throughout. Reference managers (Zotero,
  JabRef) export directly to \texttt{.bib} and will save you many hours
  over a semester.
\end{tipbox}
